\documentclass[11pt,a4paper]{article}
\usepackage[top=2.5cm,bottom=2.8cm,left=2.5cm,right=2.5cm]{geometry}
\usepackage{amsmath,amssymb,amsthm}
\usepackage{url}
\usepackage{hyperref,xcolor,enumitem,booktabs,array}
\hypersetup{colorlinks=true,linkcolor={blue!60!black},citecolor={blue!60!black},urlcolor={blue!60!black}}

\newtheorem{theorem}{Theorem}[section]
\newtheorem{proposition}[theorem]{Proposition}
\newtheorem{lemma}[theorem]{Lemma}
\newtheorem{corollary}[theorem]{Corollary}
\newtheorem{remark}[theorem]{Remark}
\theoremstyle{definition}\newtheorem{definition}[theorem]{Definition}

\newcommand{\vp}{\varphi}
\newcommand{\Lcond}{\Lambda_{\mathrm{cond}}}
\newcommand{\SU}{\mathrm{SU}}
\newcommand{\twoT}{2T}
\newcommand{\twoI}{2I}
\newcommand{\QH}{Q_H}
\newcommand{\Zthree}{\mathbb{Z}_3}
\newcommand{\Wr}{\mathrm{Wr}}
\newcommand{\Tw}{\mathrm{Tw}}
\newcommand{\Lk}{\mathrm{Lk}}
\newcommand{\nhat}{\hat{\mathbf{n}}}

\title{\textbf{The Pion-Emission Barrier of the $\QH=3$ Trefoil:
Writhe--Twist Bookkeeping, Torsion Energetics, and the Density-Feedback
Saturation of the Density-Feedback Hopfion}}
\author{Frederick Manfredi}
\date{2026}

\begin{document}
\maketitle

\begin{abstract}
Papers~XV and~XVIII establish the $\QH=3$ baryon as the $(E_6)_1$ trefoil $T_{2,3}$ and
identify its formation as a satellite-threading reconnection that topologically forces a
$\QH=1$ pion byproduct. This paper constructs that process at the field level and extracts its
framing and density-feedback energetics. (i)~A singularity-free, charge-correct \emph{knotted}
map --- the Brieskorn doublet $u=(a^2+c\,b^3)/a$ on $S^3$ --- realises the trefoil Hopfion with
$\QH=3$ (Whitehead integral, Richardson-extrapolated $Q_\infty=3.03$) and the physical virial
$V_{\max}=1.871$ matching the thin-tube value $1.875$ to $<0.3\%$, where the earlier tube-lift
and one-dimensional reductions fail. (ii)~A single-valued reaction coordinate interpolates the
$\QH=2$ Hopf link to the $\QH=3$ trefoil, realising the pion as an explicit field process: the
net Hopf charge changes by $+1$ across a reconnection barrier located at $\lambda^\ast\approx0.35$,
upgrading the topological byproduct count to a field-level statement. (iii)~A
Calugareanu--White decomposition shows the emitted pion, a $\QH=1$ \emph{unknot}, carries its
charge as pure \emph{twist} (round unknot $\Wr=0\Rightarrow\Tw=1$), while the trefoil carries
its charge as \emph{writhe}; the barrier is a twist$\leftrightarrow$writhe reconfiguration.
(iv)~The torsion energy, which cancels at leading order for the symmetric profile, remains a
small $O(\rho^2)$ curvature-gated correction ($\langle\kappa^2\tau^2\rangle/\langle\kappa^2\rangle
\lesssim4\%$) across the barrier --- it does not spike at the reconnection. (v)~The relaxed
trefoil tube is \emph{self-avoiding}: non-adjacent strands sit at a constant $2r_0$, well outside
the field-tube width, so the density-feedback saturation is a tube-\emph{wall} quantity, not an
inter-strand overlap. (vi)~The trefoil's \emph{intrinsic} self-consistent feedback is a robust,
medium-independent $\beta_*\rho\simeq0.54=8b^\ast$; the saturated $\beta_*\rho\sim15$ that
controls the baryon asymmetry (Paper~XXII) is the \emph{environment} feedback of the QCD-epoch
condensate, not this intrinsic floor. (vii)~Because $E_{\rm fb}=K_{\rm fb}J_4$ is a product the
$\QH=3$ object is a \emph{saddle}, not a minimum, and it exists as the index-1 \emph{transition
state} of the pion-emission reaction --- an interior feedback-energy barrier at
$\lambda^\ast\approx0.40$ with $\QH\!\to\!2\!\to\!3$ across it, converged in resolution ---
supplying, for the trefoil, the charge-conserving existence construction Paper~XXI
(open-problem (ii)) calls for. We state precisely what is established, what is numerically gated
by resolution, and what remains open.
\end{abstract}

\tableofcontents

\section{Scope}
\label{P23:sec:scope}
Paper~XV~\cite{Paper15} fixes the $\QH=3$ sector's profile sharpness $C_3^\ast=2.5062$, the
major radius $R_0=3$, and the minor radius $r_0=R_0/C_2^\ast=0.874$ of the trefoil
$T_{2,3}$; Paper~XVIII~\cite{Paper18} identifies the trefoil as the satellite obtained by
threading a $\QH=1$ unknot through a solid-torus neighbourhood of one component of the $\QH=2$
Hopf link, with the $(2,3)$-cable framing, and proves that Hopf-charge conservation forces a
$\QH=1$ pion byproduct. Paper~XX~\cite{Paper20} shows that the density-feedback energy of the
$\QH=3$ sector is a \emph{product} $E_{\rm fb}=K_{\rm fb}\,J_4$, so the configuration is a
\emph{saddle}, not a Derrick minimum. The present paper supplies the field-level construction
of the formation process and its energetics, to the geometric tightness of Paper~XIV's
thick-torus correction~\cite{Paper14}: a charge-correct knotted map
(Section~\ref{P23:sec:construction}), the reaction coordinate realising the pion
(Section~\ref{P23:sec:barrier}), the framing bookkeeping
(Section~\ref{P23:sec:writhetwist}), the torsion energy off the symmetric profile
(Section~\ref{P23:sec:torsion}), the self-avoidance of the relaxed tube
(Section~\ref{P23:sec:selfavoid}), and the density-feedback saturation
(Section~\ref{P23:sec:feedback}), which clarifies \emph{which} $\beta_*\rho$ enters the baryon
asymmetry of Paper~XXII~\cite{Paper22}. Numerical values quoted as resolution-gated are flagged
as such; the sector's own saddle search remains the compute-gated frontier
(Section~\ref{P23:sec:open}).

\section{A charge-correct knotted construction}
\label{P23:sec:construction}
A tube-lift of the Battye--Sutcliffe profile around the $(2,3)$ centreline (the construction of
Paper~XX) places the energy density correctly but, at accessible resolution, returns a Whitehead
charge $Q_H\approx0$: the discretised linking of the preimage fibres under-resolves. A
one-dimensional torus reduction, usable for $\QH=2$, discards the self-interaction at the
crossings and the cable winding and is invalid for the trefoil. We therefore build the field as
a projective rational map on $S^3$.

\begin{definition}[Brieskorn trefoil doublet]
\label{P23:def:brieskorn}
On $S^3\subset\mathbb{C}^2$ with $a=2(x+iy)/(1+r^2)$, $b=((1-r^2)+2iz)/(1+r^2)$,
$|a|^2+|b|^2=1$, define the normalised doublet
\begin{equation}
  (Z_1,Z_2)=\frac{(a,\;a^2+c\,b^3)}{\sqrt{|a|^2+|a^2+c\,b^3|^2}},\qquad
  \nhat=\bigl(2\,\mathrm{Re}\,\overline{Z_1}Z_2,\;2\,\mathrm{Im}\,\overline{Z_1}Z_2,\;
  |Z_1|^2-|Z_2|^2\bigr),
  \label{P23:eq:brieskorn}
\end{equation}
with $c>0$. The map $u=Z_2/Z_1=(a^2+c\,b^3)/a$ has as its $u=0$ preimage the zero set
$\{a^2+c\,b^3=0\}\cap S^3$, the link of the Brieskorn singularity, i.e.\ the $(2,3)$ torus knot.
\end{definition}

The projective normalisation removes the division singularity where $a$ or $b$ vanishes, so
$\nhat$ is everywhere well defined.

\begin{theorem}[Charge and proportion]
\label{P23:thm:charge}
The field~\eqref{P23:eq:brieskorn} has Hopf charge $\QH=3$, and its $u=0$ preimage is the knotted
trefoil (off-axis, with the two-fold toroidal modulation of $T_{2,3}$), not the axisymmetric
unknot. The modulation constant $c=3$ matches the physical trefoil virial $V_{\max}=I_0/I_2=1.875$
(Paper~XV) to better than $0.3\%$.
\end{theorem}

\begin{proof}[Numerical verification]
The Whitehead integral evaluated on~\eqref{P23:eq:brieskorn} rises with grid size along the
slow $\sim1/N$ trajectory characteristic of the validated reference map $A_{1,3}$; on a box of
half-width $L=5$ the Richardson extrapolation gives $Q_\infty=3.03$, an integer to within the
discretisation error. The $u=0$ preimage is traced to lie off the symmetry axis at
$\langle\rho\rangle\approx2.2$ with a two-fold azimuthal modulation --- the $p=2$ toroidal
winding of $T_{2,3}$ --- distinguishing it from the on-axis uniform $A_{1,3}$ unknot. With $c=3$
the virial at zero feedback is $V_{\max}=1.871$ (script \path{knotted_map.py}); pure Brieskorn
$c=1$ gives $1.80$.
\end{proof}

\section{The reaction coordinate and the pion as a field process}
\label{P23:sec:barrier}
The Brieskorn exponents label the satellite twist family of Paper~XVIII: $\{a^2+b^n=0\}$ is the
$\QH=2$ Hopf link for $n=2$, the $\QH=3$ trefoil for $n=3$, and the unfused reconnection barrier
for the intermediate state. A single-valued interpolation is therefore

\begin{proposition}[Reaction coordinate]
\label{P23:prop:reaction}
The family
\begin{equation}
  u(\lambda)=\frac{a^2+(1-\lambda)\,b^2+\lambda\,c\,b^3}{a},\qquad\lambda\in[0,1],\quad c=3,
  \label{P23:eq:reaction}
\end{equation}
interpolates the Hopf link ($\lambda=0$, $\QH\to2$) to the trefoil ($\lambda=1$, $\QH\to3$). The
phase of the Brieskorn numerator winds once with the toroidal angle only when the $b^3$ term
dominates; that onset of winding is the reconnection at which the two Hopf-link components fuse
into a single knot.
\end{proposition}

\begin{theorem}[Pion emission as a field process]
\label{P23:thm:pion}
Along~\eqref{P23:eq:reaction} the Hopf charge changes by $\Delta\QH=+1$ between the link and the
knot. The $+1$ is the $\QH=1$ pion of Paper~XVIII, realised not as a bookkeeping byproduct but as
an explicit change of the field's topological charge across the barrier.
\end{theorem}

\begin{proof}[Numerical verification]
On a box $L=5$ at $N=128$ the gate returns $\QH=1.93$ at $\lambda=0$ (Hopf link) and $2.86$ at
$\lambda=1$ (trefoil); finite resolution under-counts $2$ and $3$ equally, but the difference
$+0.93\approx+1$ is robust (script \path{barrier_config.py}). The reconnection --- the steepest
$d\QH/d\lambda$, coincident with the peak of the self-interaction energy $E_{\rm fb}$ --- locates
the barrier at $\lambda^\ast\approx0.35$.
\end{proof}

\begin{remark}[The barrier configuration]
\label{P23:rem:barrier}
The configuration $\nhat(\lambda^\ast)$ is the half-threaded state ``just before full absorption /
just after emission begins''. It is the configuration on which the density-feedback saturation of
Section~\ref{P23:sec:feedback} and the baryon asymmetry of Paper~XXII are to be evaluated.
\end{remark}

\section{Writhe--twist bookkeeping: the pion is a twist quantum}
\label{P23:sec:writhetwist}
For a framed knotted tube the Hopf charge is the self-linking of the framed centreline, which by
the Calugareanu--White theorem splits as $\Lk=\Wr+\Tw$ (writhe plus framing twist).

\begin{theorem}[The pion carries twist]
\label{P23:thm:calugareanu}
The emitted $\QH=1$ pion is the unknot. In its ground (round) state $\Wr=0$, so by
$\Lk=\Wr+\Tw$ its unit of Hopf charge is pure \emph{twist}, $\Tw=1$. On emission it removes one
twist quantum; on absorption it supplies one. The trefoil, by contrast, carries its $\QH=3$
predominantly as writhe.
\end{theorem}

\begin{proposition}[Writhe--twist trend across the barrier]
\label{P23:prop:wrtw}
The Gauss self-writhe and linking of the $u=0$ preimage centreline, computed along the reaction
coordinate (script \path{torsion_framing.py}), give: at $\lambda=0$ the charge is pure mutual
linking ($\Wr_{\rm self}=0$, $2\Lk=2.000=\QH$, $\Tw=0$); the centreline topology reconnects
(link $\to$ knot) at $\lambda\approx0.3$, confirming the barrier of
Theorem~\ref{P23:thm:pion} by an independent diagnostic; at $\lambda=1$ the single trefoil is
writhe-dominated, $\Wr_{\rm self}=3.77$ on this embedding, with a compensating framing twist
$\Tw=\QH-\Wr=-0.77$. The reconnection thus converts mutual linking into self-writhe while the net
charge rises by exactly $+1$, carried off by the unknotted pion as a twist quantum.
\end{proposition}

The writhe value is embedding-dependent, hence so is the specific residual twist; the
frame-independent content is (i)~the link end is pure linking, (ii)~the reconnection at
$\lambda\approx0.3$, and (iii)~the pion, being the unknot, carries its charge as twist.

\section{Torsion energetics off the symmetric profile}
\label{P23:sec:torsion}
Paper~XV proves that for the $\chi$-independent Battye--Sutcliffe profile the Frenet torsion acts
as a rigid rotation of the tube cross-section and makes no contribution to the energy at any
order. That cancellation is a leading-order statement; the barrier configuration is not symmetric,
so the question is whether torsion remains negligible there.

\begin{proposition}[Torsion stays perturbative at the barrier]
\label{P23:prop:torsion}
The surviving torsion correction to the $\chi$-averaged gradient energy is curvature-gated,
$\langle\delta|\nabla f|^2\rangle=\tfrac12 f_1(\rho)^2\kappa(s)^2\tau(s)^2\rho^2$ (Paper~XV); the
relevant arc-length average $\langle\kappa^2\tau^2\rangle/\langle\kappa^2\rangle$ is $4.4\%$ for
the relaxed symmetric trefoil ($\langle\kappa^2\rangle=0.120$,
$\langle\kappa^2\tau^2\rangle=0.0053$). Tracked along the reaction coordinate --- where the
cancellation does not apply --- this ratio remains in the $0$--$3\%$ band \emph{including} the
barrier; it does not spike at the reconnection. The torsion/framing energy is therefore a small
$O(\rho^2)$ correction throughout, not a leading term.
\end{proposition}

\begin{proof}[Numerical verification]
A numerical Frenet evaluation of the analytic $(2,3)$ centreline reproduces the Paper~XV values
$\langle\kappa^2\rangle=0.1200$, $\langle\kappa^2\tau^2\rangle=0.0053$ and arc length
$L_3=41.26$ exactly (script \path{torsion_framing.py}); applied to the reaction centreline,
rescaled to $R_0=3$, the ratio is $0\%$ at the Hopf link (planar circles, $\tau=0$), $3.2\%$ at
the barrier $\lambda=0.3$, and small at the trefoil, converged over grid refinement. The
Brieskorn embedding is geometrically smoother than the physical trefoil, so the robust statement
is the \emph{shape} (no spike, bracketed $0$--$4\%$), not the absolute percentage.
\end{proof}

\section{The relaxed tube is self-avoiding}
\label{P23:sec:selfavoid}
Where the feedback saturates is decided by where the field density is largest. For the trefoil
this is \emph{not} at inter-strand crossings.

\begin{proposition}[Self-avoidance]
\label{P23:prop:selfavoid}
For the relaxed $(2,3)$ tube ($R_0=3$, $r_0=0.874$) the minimal distance between non-adjacent
strands is the constant $|G(t)-G(t+\pi)|=2r_0=1.748$ for all $t$, well outside the field-tube
width $2/C_3^\ast=0.80$. The relaxed knot is therefore self-avoiding: the field tubes do not
overlap, and the density-feedback saturation is a single-tube \emph{wall} quantity, roughly
uniform along the knot, rather than a crossing-overlap quantity. The ``self-interaction at the
crossings'' of Paper~XV is the \emph{curvature} peaking there
($\kappa_{\max}=0.399$, already in $\langle\kappa^2\rangle$), not a tube overlap.
\end{proposition}

Inter-strand overlap --- and hence an overlap enhancement of the feedback --- occurs only
transiently, during the threading of the barrier configuration
(Remark~\ref{P23:rem:barrier}), not in the relaxed endpoint.

\begin{remark}[Single knotted circle, two cable strands]
\label{P23:rem:cable}
The trefoil preimage is a \emph{single knotted circle} (Paper~XVIII~\cite{Paper18}), not two
components; the two strands at separation $2r_0$ are the two passes of the $(2,3)$-cable through
each toroidal section, at the inner radius $R_0-r_0\cos3t$ and the outer radius
$R_0+r_0\cos3t$. The construction here is therefore single-tube, with the cable's two strands
captured by the two nearest centreline passes. Baryon formation is the \emph{fusion} of these
two strands: the two-component Hopf link $T_{2,2}$ ($\QH=2$) knots into the single-component
trefoil $T_{2,3}$ ($\QH=3$) with the $\QH=1$ pion left over
(Section~\ref{P23:sec:barrier}), and the transient inter-strand overlap of that fusion ---
absent in the relaxed, self-avoiding endpoint --- is where the feedback saturation is enhanced.
\end{remark}

\section{Density-feedback saturation: intrinsic floor versus environment}
\label{P23:sec:feedback}
The density-feedback denominator $1+\beta_*\kappa$, with $\kappa=|\nabla\nhat|^2$ the
Faddeev--Niemi gradient density (Paper~VII~\cite{Paper7}), controls both the directional
suppression $S_{\rm eff}$ and the saddle condition of Paper~XX.

\begin{proposition}[Intrinsic self-consistent feedback is a robust floor]
\label{P23:prop:floor}
The trefoil's \emph{intrinsic} self-consistent feedback argument is scale- and
medium-independent, $\beta_*\kappa_{\rm core}=8b^\ast\simeq0.54$. The Battye--Sutcliffe core
kernel is bounded, $(f_0')^2+\sin^2\!f_0/\rho^2=8/C^{\ast2}$, and self-consistency
$J_{fb}/J_a=\vp$ fixes $b^\ast=\beta_*/C^{\ast2}=0.0672$ (Paper~IV~\cite{Paper4}); the product
$8b^\ast$ is independent of $C^\ast$, so no sharpening or compression of the knot raises it.
\end{proposition}

\begin{proof}[Numerical verification]
On the physical-tube field the self-consistency $V=J_{2\rm iso,fb}/J_{2a}=\vp$ is solved by a
coupling $\beta_{\rm sc}$ at which the feedback argument at the core is $0.45$--$0.50$ (script
\path{feedback_energy_3d.py}), matching the analytic $8b^\ast=0.54$. Adding a uniform ambient
background to the self-consistency does not raise it: the coupling adjusts to hold $V=\vp$ and the
feedback argument stays $\simeq0.3$--$0.5$ for any background (script \path{knot_in_medium.py}).
\end{proof}

\begin{remark}[The baryon-asymmetry feedback is the environment value]
\label{P23:rem:environment}
The intrinsic floor of Proposition~\ref{P23:prop:floor} is \emph{not} the $\beta_*\rho\sim15$ that
sets the baryon asymmetry (Paper~XXII). The feedback coupling is a chameleon
(Paper~VII~\cite{Paper7}): $\beta_*\rho$ runs from $\vp$ at the vacuum to $\sim10^6$ at
Solar-System density, so the relevant quantity at baryogenesis is the \emph{environment}
feedback of the dense QCD-epoch condensate in which the trefoil forms. That value is bracketed by
the first-order-transition window and freeze-out-selected at $T_c$ (the baryon, a confined-phase
knot, imprints the asymmetry as it crystallises out of the deconfined soup), landing at
$\beta_*\rho\sim15$ for the observed $\eta$. The intrinsic trefoil feedback ($\simeq0.54$) is a
subdominant addition to this environment value; the two must not be conflated.
\end{remark}

\section{Existence of the $\QH=3$ configuration as a pion-emission transition state}
\label{P23:sec:existence}
The density-feedback energy $E_{\rm fb}=K_{\rm fb}J_4$ is a \emph{product}, so the $\QH=3$ Hopfion
is a \emph{saddle}, not a Derrick minimum (Paper~XX~\cite{Paper20}); a static energy
minimisation therefore does not locate it. Its natural realisation is the index-1 saddle --- the
\emph{transition state} --- of the pion-emission reaction of Section~\ref{P23:sec:barrier}. A
static search is moreover \emph{degenerate} for a sharper reason, which also fixes how the feedback
coupling must be treated:

\begin{lemma}[Scale--coupling degeneracy of the feedback energy]
\label{P23:lem:scaling}
Under the dilation $\mathbf{n}(\mathbf{x})\to\mathbf{n}(\mathbf{x}/\lambda)$ the feedback integrals
scale as $J_{2a}\to\lambda\,J_{2a}$, $J_4\to\lambda^{-1}J_4$, and
$J_{2\mathrm{iso,fb}}(\beta)\to\lambda\,J_{2\mathrm{iso,fb}}(\beta/\lambda^2)$, so the virial ratio
obeys $V(\lambda,\beta)=V(\beta/\lambda^2)$. A dilation is therefore equivalent to a rescaling of
the feedback coupling, $\beta\to\beta/\lambda^2$.
\end{lemma}

\begin{proof}
With grid spacing $h\to\lambda h$ representing the dilation, $|\nabla\mathbf{n}|^2\to\lambda^{-2}
|\nabla\mathbf{n}|^2$ and $dV\to\lambda^3 dV$, while the pointwise $\sin^4\!\theta$ and
$\mathbf{n}$ are unchanged; substituting into $J_{2a}=\int\sin^4\!\theta\,|\nabla\mathbf{n}|^2\,dV$,
$J_4=\int F^2\,dV$ (with $F$ bilinear in gradients, $F\to\lambda^{-2}F$), and
$J_{2\mathrm{iso,fb}}(\beta)=\int|\nabla\mathbf{n}|^2/(1+\beta|\nabla\mathbf{n}|^2)\,dV$ gives the
stated factors, whence $V=J_{2\mathrm{iso,fb}}/J_{2a}\to J_{2\mathrm{iso,fb}}(\beta/\lambda^2)/
J_{2a}=V(\beta/\lambda^2)$. (Verified numerically: $V(\lambda h,\beta)=V(h,\beta/\lambda^2)$ to
machine precision, \path{existence_saddle.py}.) \qed
\end{proof}

\noindent This is the density-feedback form of the Derrick/dilation instability that makes the
$\QH=2$ Hopfion a saddle (Paper~I~\cite{Paper1}): the self-consistency $V=\vp$ fixes the scale once $\beta$ is
given, so a stationary-point search over scale \emph{and} $\beta$ together drives the score
$|V-\vp|+|s_{\rm opt}-1|$ to zero for \emph{any} shape and certifies nothing. The physical coupling
$\beta$ is held fixed --- calibrated in-engine and shown to be a saddle-structure quantity
(Section~\ref{P23:sec:feedback}) --- and only the shape is varied; the existence statement below
rests instead on the reaction coordinate.

\begin{proposition}[Existence as a transition state]
\label{P23:prop:existence}
Along the reaction coordinate~\eqref{P23:eq:reaction} the feedback energy $E_{\rm fb}(\lambda)$ has
an interior maximum (a barrier) at $\lambda^\ast\approx0.40$, across which the Hopf charge passes
from $\QH\!\to\!2$ (Hopf link) to $\QH\!\to\!3$ (trefoil) with net change $+1$ (the pion). This
interior maximum, separating the two charge sectors, is an index-1 saddle of $E_{\rm fb}$: the
$\QH=3$ configuration exists as the pion-emission transition state --- a \emph{saddle}, not the
energy minimiser a static picture would seek. This supplies, for the trefoil, the
charge-conserving construction and existence demonstration that Paper~XXI~\cite{Paper21}
(open-problem (ii)) identifies as required --- there noted to fail under charge-non-conserving
free-flow relaxation, ``even for the known trefoil.'' The analogous construction for the D-type
sector of Paper~XXI remains to be carried out.
\end{proposition}

\begin{proof}[Numerical verification]
Evaluating $E_{\rm fb}$ on~\eqref{P23:eq:reaction} (script \path{existence_saddle.py}, mode
\texttt{-{}-reaction}) at $N=128,160,192$: the barrier sits at $\lambda^\ast=0.40$ at every
resolution (its location is converged), with $E_{\rm fb}$ rising from the Hopf-link value to the
peak and descending on the trefoil side; the charge change is $\Delta\QH=+0.93,\,+0.96,\,+0.97$,
climbing to the pion unit $+1$, with endpoint charges $1.93,1.95,1.96\!\to\!2$ and
$2.86,2.91,2.93\!\to\!3$ (finite-$N$ under-counts both equally). By the mountain-pass argument an
index-1 saddle lies between the two basins; the reaction coordinate is one path over it, so the
saddle energy is at most the computed barrier.
\end{proof}

\begin{remark}[Scope of the existence statement]
\label{P23:rem:existence_scope}
The transition-state saddle is the robust statement. A \emph{static} self-consistent configuration
($V=\vp$ and $s_{\rm opt}=1$) is \emph{not} found within the charge-certified ansatz family: the
Brieskorn family has $s_{\rm opt}\sim4$--$9$, far from unity, while the energy-localised tube-lift
is charge-broken ($\QH\approx0$). This dichotomy --- charge-correct or energy-saddle-like, not both
at accessible resolution --- is the trefoil resolution wall, and is consistent with the object
being a saddle rather than a minimum. A post-barrier trefoil-side basin (near $\lambda\approx0.8$)
emerges at $N\ge160$ but is shallow and resolution-dependent; the firm result is the
transition-state saddle and the converged $\lambda^\ast$.
\end{remark}

\section{Status and open problems}
\label{P23:sec:open}
The topological and framing results of Sections~\ref{P23:sec:construction}--%
\ref{P23:sec:selfavoid} are established: the charge-correct knotted map, the pion as a field
process, the twist character of the pion, the perturbative torsion, and the self-avoidance; and
the $\QH=3$ configuration exists as a pion-emission transition state
(Section~\ref{P23:sec:existence}). Three quantities remain open.
\begin{enumerate}[leftmargin=2em]
\item \textbf{The transition-state energy in physical units, and the first-order window location.}
Proposition~\ref{P23:prop:existence} establishes existence as a transition state; what remains is
its \emph{physical-units} energy --- the sector's own density-feedback (formation) energy, which
fixes the physical location of the first-order window of Paper~XXII. This is \emph{not} a static
relaxation (the trefoil, being transient, has no isolated minimum to relax to); it is a
formation-energy evaluation of the transition state, gated by the tube resolution wall (Paper~XX),
and remains open.
\item \textbf{Absolute writhe and torsion.} The self-writhe ($3.77$) and the residual twist
($-0.77$) are embedding-dependent on the Brieskorn centreline; their physical values require the
relaxed centreline geometry.
\item \textbf{The barrier energetics.} The density-feedback energy evaluated on the relaxed
half-threaded barrier configuration --- the transient inter-strand overlap of
Section~\ref{P23:sec:selfavoid} --- is the remaining piece connecting the construction here to the
baryon-asymmetry magnitude of Paper~XXII; it too awaits the relaxed saddle.
\end{enumerate}
What is \emph{not} open is the identification of the quantities: the baryon-asymmetry feedback is
the environment value, not the intrinsic floor (Remark~\ref{P23:rem:environment}), and the pion is
a twist quantum (Theorem~\ref{P23:thm:calugareanu}).

\begin{thebibliography}{99}
\bibitem{Paper1} F.~Manfredi, \textit{The Density-Feedback Faddeev--Niemi Hopfion:
Exact Algebraic Predictions for Standard-Model Parameters}, Zenodo (2026).
\href{https://doi.org/10.5281/zenodo.19342027}{doi:10.5281/zenodo.19342027}
\bibitem{Paper4} F.~Manfredi, \textit{The Hopf Spoke, WZW Fermion Mass Renormalization,
and Three- and Four-Term Formulas for the Fine Structure Constant}, Zenodo (2026).
\href{https://doi.org/10.5281/zenodo.19504729}{doi:10.5281/zenodo.19504729}
\bibitem{Paper7} F.~Manfredi, \textit{Emergent Gravity, Inflation, Dark Energy, and the Weak
Equivalence Principle from the Density-Feedback Hopfion Condensate}, Zenodo (2026).
\href{https://doi.org/10.5281/zenodo.19646953}{doi:10.5281/zenodo.19646953}
\bibitem{Paper14} F.~Manfredi, \textit{The Thick-Torus Profile Correction to the
Density-Feedback Hopfion}, Zenodo (2026).
\href{https://doi.org/10.5281/zenodo.20479790}{doi:10.5281/zenodo.20479790}
\bibitem{Paper15} F.~Manfredi, \textit{The $Q_H=3$ Sector of the Density-Feedback Hopfion},
Zenodo (2026). \href{https://doi.org/10.5281/zenodo.20691001}{doi:10.5281/zenodo.20691001}
\bibitem{Paper18} F.~Manfredi, \textit{Preimage Topology and Confinement: The $Q_H=3$ Sector
of the Density-Feedback Hopfion}, Zenodo (2026).
\href{https://doi.org/10.5281/zenodo.21047750}{doi:10.5281/zenodo.21047750}
\bibitem{Paper20} F.~Manfredi, \textit{The Dynamical Layer of the Quark Sector: $E_6$-Native
Generations, the Isospin Mass-Scale, and Static Residuals as Strong Dynamics}, Zenodo (2026).
\href{https://doi.org/10.5281/zenodo.23047960}{doi:10.5281/zenodo.23047960}
\bibitem{Paper21} F.~Manfredi, \textit{The Beyond--Standard-Model D-type Sector of the
Density-Feedback Hopfion: A Sterile, Dirac Dark-Matter Candidate}, Zenodo (2026).
\href{https://doi.org/10.5281/zenodo.23319530}{doi:10.5281/zenodo.23319530}
\bibitem{Paper22} F.~Manfredi, \textit{Dark Genesis: Co-genesis of Baryons and a Sterile Dirac
Dark Sector at the First-Order Confinement Transition of the Density-Feedback Hopfion}, Zenodo
(2026). \href{https://doi.org/10.5281/zenodo.23319779}{doi:10.5281/zenodo.23319779}
\end{thebibliography}

\end{document}
